\section{Validity and scope}
\label{sec:validity}

The construction has two distinct claims.  First, every admissible parameter
vector defines an arbitrage-free slice.  Second, increasing polynomial order
fills out a specified class of laws.  The first is exact at every order; the
second determines the scope of the coordinates.

\subsection{One-expiry validity}

\begin{proposition}[One-expiry no-arbitrage]
\label{prop:noarb}
Let $\theta\in\Theta_N$ and price calls from the LQD slice
\eqref{eq:lqd}.  Then:
\begin{enumerate}[label=(\roman*)]
\item $X$ has a strictly positive continuous density on $\R$, and
$Y=e^X$ has a strictly positive density on $(0,\infty)$ with $\E[Y]=1$;
\item as a function of strike $y$, the call is strictly decreasing and
strictly convex, satisfies $(1-y)^+\le c(y)\le1$, has
$c'(y)=-\Prob(Y>y)\in(-1,0)$, and tends to zero as $y\to\infty$;
\item every nondegenerate finite-width butterfly has positive price, and
put--call parity holds identically;
\item as a function of log-strike $k$, the call is concave sufficiently far
in the money.
\end{enumerate}
\end{proposition}

\begin{proof}
Equation \eqref{eq:density} gives a positive continuous density; the shift of
\cref{prop:shift} gives unit mean.  Rank integration gives
\[
c(y)=\int_0^1(e^{Q(u)}-y)^+\dd u.
\]
Each integrand is decreasing and convex in $y$, and differentiation yields
$c'(y)=-\Prob(Y>y)$.  Positivity of the density makes both monotonicity and
convexity strict.  Jensen gives $(1-y)^+\le c(y)$, while
$c(y)\le\E[Y]=1$ and dominated convergence gives the right boundary.
Butterfly positivity and parity follow immediately.  Finally,
\[
\frac{\partial^2c}{\partial k^2}
=e^k\{f_X(k)-\Prob(Y>e^k)\},
\]
which is negative as $k\to-\infty$ because $f_X(k)\to0$ and the probability
tends to one.
\end{proof}

The proposition is stronger than a dense-grid test.  It covers every strike,
including strikes between grid nodes and beyond the displayed smile, because
the call is an expectation under a positive law.  Calibration may move the
coefficient vector arbitrarily inside $\Theta_N$ without approaching a hidden
butterfly boundary.  The numerical implementation must still preserve this
continuous geometry; that separate question is treated by the certificates
of \cref{sec:certificates}.

\begin{figure}[htbp]
\centering
\paperfig{0.92\textwidth}{fig_butterfly.pdf}
\caption{Validity on the fitted SPY December 2026 slice.  Left: the call
curve is convex in normalized strike and lies below its secant chord.  Right:
finite butterflies, divided by squared half-width, converge to the density.
The geometry is produced by the law, not imposed during calibration.}
\label{fig:butterfly}
\end{figure}

\subsection{Why the wall cannot be removed}

\begin{proposition}[Finite-forward wall]
\label{prop:wall}
Within the LQD family the martingale normalizer exists if and only if
$\lambda_+<1$.  More generally, if a law satisfies
$\Prob(X>x)\ge\kappa e^{-x}$ for some $\kappa>0$ and all sufficiently large
$x$, then $\E[e^X]=\infty$.
\end{proposition}

\begin{proof}
The first statement is \cref{prop:shift}.  For the second, at any large
$x_0$,
\[
\E[e^X]\ge\int_{e^{x_0}}^\infty\Prob(X>\log t)\,\dd t
\ge\kappa\int_{e^{x_0}}^\infty\frac{\dd t}{t}=\infty.
\]
\end{proof}

The boundary is therefore probabilistic, not parametric.  A right log-return
tail with scale at least one cannot support a finite forward.  There is no
left counterpart because $e^X$ is bounded on the negative half-line.

\subsection{Coverage of the coordinates}

The natural target class can be described without reference to polynomials.
For a differentiable quantile $Q^*$ define
\[
g^*(u)=\log Q^{*\prime}(u)+\log u+\log(1-u).
\]
Continuity of $g^*$ on $[0,1]$ is exactly the condition that the law have a
positive continuous density and exponential tails with finite nonzero endpoint
scales.

\begin{proposition}[Universality, with exclusions]
\label{prop:universal}
Let $X^*$ be a mean-one law for which $g^*$ extends continuously to
$[0,1]$, with $e^{g^*(1)}<1$.  Then:
\begin{enumerate}[label=(\roman*)]
\item There exist admissible LQD slices with $g_N\to g^*$ uniformly.  Their
quantiles converge uniformly on compact subsets of $(0,1)$, their call curves
converge uniformly over all strikes, and their gross-return laws converge in
Wasserstein-$p$ distance for every $1\le p<e^{-g^*(1)}$.
\item No finite-order LQD slice has an atom, a zero-density interval, bounded
support, mass at zero, or Gaussian-or-faster tails.  These exclusions persist
under uniformly convergent log-speeds.  More precisely, let admissible $g_N$
converge uniformly to $g_\infty$.  If $g_\infty(1)<0$, the limiting law remains
in the class of part (i).  If $g_\infty(1)=0$, weak convergence can fail because
mass escapes to $-\infty$; if a weak limit nevertheless exists, it has full
support and a strictly positive continuous density, but sits at the
finite-forward boundary.  Its moments of order $r>1$ diverge, and its mean may
be smaller than one.
\end{enumerate}
\end{proposition}

Here the one-dimensional Wasserstein metric has the concrete quantile form
\[
W_p(\mu,\nu)^p=\int_0^1|Q_\mu(u)-Q_\nu(u)|^p\,\dd u.
\]
The proof is in \cref{app:proofs}.  Its approximation step is Weierstrass:
polynomials approximate the continuous $g^*$ uniformly.  Since
$g^*(1)<0$, sufficiently accurate approximants remain inside the open
admissible set.  The tail margin supplies a common integrable envelope for
normalizers, calls, and moments.

Uniform convergence of $g$ is the correct strength for this result.  It
controls the transport speed at every rank, including both endpoints.  Local
approximation on the quoted ranks alone would say nothing about normalization
or tail moments.  Once a strict margin $g^*(1)<0$ is fixed, the same exponential
envelope controls all approximating normalizers and supplies uniform
integrability of gross returns.  This is what upgrades pointwise shape
approximation to uniform convergence of call prices.

The exclusions follow from the same geometry.  A positive continuous
transport speed makes the quantile strictly increasing and unbounded.  An
atom would require a quantile jump; a gap, a flat interval; bounded support, a
finite endpoint; and a Gaussian tail, $g(u)\to-\infty$ at an endpoint.  None is
available to a finite polynomial, nor to a uniform limit of bounded log-speeds.

Multimodality is not excluded.  By \eqref{eq:reciprocal}, density modes are
created where the transport slows and are separated where it accelerates.
A positive speed can vary many times without ever becoming zero.  Thus the
family may represent a smooth double-hump density, but not a literal gap of
zero density between the humps.  Basis order determines how sharply the speed
can vary, as the synthetic example in \cref{sec:examples} demonstrates.

Thus the LQD family is not a coordinate system for every probability law.  It
is a coordinate system for a flexible exponential-tail class.  That class is
large enough to approximate smooth unimodal and multimodal smile laws, while
its boundary and exclusions remain explicit.
